\section{Quantitative analysis}

Quantitative analysis turns process models into numbers: throughput, waiting times, utilisation, cost.\footnote{M.~Dumas, M.~La~Rosa, J.~Mendling, H.~A.~Reijers, \emph{Fundamentals of Business Process Management}, Ch.~7; course material at \url{https://courses.cs.ut.ee/2014/bpm/}.} Three concerns are easy to conflate. \textbf{Validation} (model versus reality) and \textbf{verification} (the soundness-style questions of the earlier chapters: can it deadlock, can every case terminate?) belonged to those chapters. \textbf{Performance analysis} asks quantitative ones: how many cases per hour, what average flow time, how many extra resources to reach a given throughput.

\subsection{Performance measures and KPIs}

\begin{definitionbox}{Performance dimensions, measures, KPI}
Companies want processes \emph{faster}, \emph{cheaper}, \emph{better}: the three performance dimensions \textbf{time}, \textbf{finance}, \textbf{quality}. A \textbf{process performance measure} is a quantity that can be unambiguously determined for a given business process; the abstract dimensions are refined into concrete \textbf{Key Performance Indicators} (KPI): scalar quantities computable from process executions and usable as optimisation targets.
\end{definitionbox}

\begin{definitionbox}{Cycle, processing, waiting time}
The time dimension decomposes into three measures: \textbf{cycle time}, the time needed to handle one case from start to end; \textbf{processing time} (or \emph{service time}), the time resources spend actually handling the case; \textbf{waiting time}, the time the case spends idle, including queueing caused by resource unavailability. Cycle time is the sum of the other two.
\end{definitionbox}
The split matters because the two parts respond to different interventions: processing time is hard to compress without changing what a task does, waiting time often yields to better scheduling or extra resources. A cycle-time goal can take three shapes: reduce the \emph{average}, reduce the \emph{maximum}, or meet a deadline \emph{negotiated with the customer} (a service-level agreement), and they drive different redesigns: averages hide outliers, a maximum or a contract forces attention to the worst cases.

The finance dimension covers \emph{cost}, \emph{turnover}, \emph{yield} and \emph{revenue}; a yield increase moves profit as much as a cost decrease, but redesign is most often framed in terms of cost. Costs split into \textbf{fixed} (overhead unaffected by processing intensity: infrastructure, maintenance), \textbf{variable} (positively correlated with sales volume, purchases, hires) and \textbf{operational} (directly related to the output of the process, labour cost above all). Redesign typically targets operational cost. Automating a task cuts labour cost but introduces \emph{incidental} cost (developing the replacing application) and fixed \emph{maintenance} cost for it: automation pays only when case volume amortises both over the saved labour.

Quality has two viewpoints. \textbf{External quality} is the client's satisfaction with the product or with how the process ran; key factors are the amount, relevance, quality and timeliness of the information the client receives as the process progresses. \textbf{Internal quality} is the participants': level of control over the work, variation experienced, challenges faced. External quality is often measured through time (average cycle time, percentage of missed deadlines), and by convention any time-based measure is classified under the \emph{time} dimension.

Each measure is paired with an \textbf{aggregation function} turning a population of executions into one scalar: count, average, variance, minimum, maximum. ``Average delivery cost per item'' pairs cost-per-item with the average; ``maximal time of delivery'' pairs delivery time with the maximum. Deriving a KPI from a management goal takes three steps: \emph{formulate} the performance objective (the desirable state of the process), \emph{identify} the relevant dimension and aggregation function and derive the KPI, \emph{set} a target value.
\begin{examplebox}{Restaurant KPI}
A restaurant is losing customers to poor service, and the management team focusses on meal delivery. Customer interviews give two anchors: half the customers want their meal within $15$ minutes, and all agree that $30$ minutes or more is unacceptable. The same goal supports three different KPI:
\begin{center}
\begin{tabular}{lllc}
\toprule
 & \textbf{Objective} & \textbf{KPI (time dimension)} & \textbf{Target}\\
\midrule
v1 & served within $30'$ & $ST_{30}$ = \% of customers served in $<30'$ & $\geq 97\%$\\
v2 & served in about $15'$ & $ST_{15}$ = \% of customers served in $<15'$ & $\geq 85\%$\\
v3 & served in about $15'$ & $AMDT$ = average meal delivery time & $\leq 15'$\\
\bottomrule
\end{tabular}
\end{center}
v1 guards the tail (almost nobody waits unacceptably long, nothing promised about typical service); v2 tightens the threshold but relaxes the target, tolerating a long tail; v3 collapses the distribution into an average.
\end{examplebox}
The canonical measures in practice: for time, cycle time and waiting time (time in non-value-added tasks); for cost, cost per execution and resource utilisation; for quality, error rates and missed promises.

\subsection{Cycle time analysis}

\begin{definitionbox}{Flow analysis}
\textbf{Flow analysis} estimates the overall performance of a process from the performance of its individual activities: cycle time of the process from the cycle times of the activities, average cost from costs per execution, error rate from activity error rates. It trades realism for tractability, assuming that activities are independent and that the model's structure dictates how their numbers compose.
\end{definitionbox}
\begin{definitionbox}{Cycle time analysis}
The \textbf{cycle time} of a case is the difference between its end time and its start time; \textbf{cycle time analysis} computes the average cycle time of a process or fragment from the cycle times of its activities, assuming an average \emph{activity time} (waiting $+$ processing) is known for every activity.
\end{definitionbox}
Notation: $CT$ is the average cycle time of the process; $CT_A$, $CT_B, \ldots$ those of activities $A$, $B, \ldots$ In diagrams the activity time appears in parentheses under the name, $A\,(10)$. Composition proceeds by \emph{flow patterns}: sequence, alternative paths (XOR-block), parallel paths (AND-block), rework loops.

\paragraph{Sequence.} Time composes additively along a sequence: a case reaching $B$ has already spent $CT_A$ in $A$. For $A\,(10)$ followed by $B\,(20)$, $CT = 10 + 20 = 30$. In general
\[CT = \sum_{i=1}^{n} CT_i.\]

\paragraph{XOR-block and AND-block.} The fragment between an XOR-split and its matching XOR-join is an \textbf{XOR-block}; between an AND-split and its AND-join, an \textbf{AND-block}.
\begin{center}
\begin{tikzpicture}[node distance=0.3cm and 0.4cm]
  \node[bpmn event] (s) {};
  \node[bpmn task, right=of s] (A) {$A\,(10)$};
  \node[bpmn gateway, right=of A] (g1) {$\times$};
  \node[bpmn task, above right=of g1] (B) {$B\,(20)$};
  \node[bpmn task, below right=of g1] (C) {$C\,(15)$};
  \node[bpmn gateway, below right=of B] (g2) {$\times$};
  \node[bpmn event, right=of g2] (e) {};
  \draw[bpmn flow] (s) -- (A);
  \draw[bpmn flow] (A) -- (g1);
  \draw[bpmn flow] (g1) -- (B);
  \draw[bpmn flow] (g1) -- (C);
  \draw[bpmn flow] (B) -- (g2);
  \draw[bpmn flow] (C) -- (g2);
  \draw[bpmn flow] (g2) -- (e);
\end{tikzpicture}\hspace{0.8cm}
\begin{tikzpicture}[node distance=0.3cm and 0.4cm]
  \node[bpmn event] (s) {};
  \node[bpmn task, right=of s] (A) {$A\,(10)$};
  \node[bpmn gateway, right=of A] (g1) {$+$};
  \node[bpmn task, above right=of g1] (B) {$B\,(20)$};
  \node[bpmn task, below right=of g1] (C) {$C\,(15)$};
  \node[bpmn gateway, below right=of B] (g2) {$+$};
  \node[bpmn event, right=of g2] (e) {};
  \draw[bpmn flow] (s) -- (A);
  \draw[bpmn flow] (A) -- (g1);
  \draw[bpmn flow] (g1) -- (B);
  \draw[bpmn flow] (g1) -- (C);
  \draw[bpmn flow] (B) -- (g2);
  \draw[bpmn flow] (C) -- (g2);
  \draw[bpmn flow] (g2) -- (e);
\end{tikzpicture}
\end{center}
In the XOR-block each case takes exactly one branch, so its cycle time is $30$ (via $B$) or $25$ (via $C$); where the \emph{average} sits depends on how often each branch is taken. The \textbf{branching probability} $p_i$ of branch $i$ is the frequency with which it is taken, with $\sum_i p_i = 1$, estimated from historical data or domain knowledge: without the $p_i$ the cycle time of an XOR-block is undefined. The rule is the weighted average
\[CT = \sum_{i=1}^{n} p_i \cdot CT_i,\]
giving $0.5 \cdot 30 + 0.5 \cdot 25 = 27.5$ for a $50/50$ split and $0.9 \cdot 30 + 0.1 \cdot 25 = 29.5$ for $90/10$. In the AND-block both branches run concurrently, so summing all three activity times ($45$) overestimates: the AND-join waits for the \emph{slowest} branch, $CT = 10 + \max\{20, 15\} = 30$. The rule is
\[CT = \max_i\, CT_i,\]
and the synchronisation cost is the time the faster branch wastes waiting.

\begin{examplebox}{Loan application}
\begin{center}
\resizebox{\linewidth}{!}{%
\begin{tikzpicture}[every node/.append style={font=\scriptsize\sffamily}]
  \node[bpmn event] (start) at (0,0) {};
  \node[bpmn task, align=center] (cc) at (1.7,0) {Check\\completeness\\(1 d)};
  \node[bpmn gateway] (as) at (3.4,0) {$+$};
  \node[bpmn task, align=center] (ch) at (5.3,0.9) {Check credit\\history (1 d)};
  \node[bpmn task, align=center] (ci) at (5.3,-0.9) {Check income\\sources (3 d)};
  \node[bpmn gateway] (aj) at (7.2,0) {$+$};
  \node[bpmn task, align=center] (aa) at (8.9,0) {Assess\\application\\(3 d)};
  \node[bpmn gateway] (xs) at (10.6,0) {$\times$};
  \node[bpmn task, align=center] (mo) at (12.6,0.9) {Make credit\\offer (1 d)};
  \node[bpmn task, align=center] (nr) at (12.6,-0.9) {Notify\\rejection (2 d)};
  \node[bpmn gateway] (xj) at (14.5,0) {$\times$};
  \node[bpmn event, draw=red] (end) at (15.5,0) {};
  \draw[bpmn flow] (start) -- (cc);
  \draw[bpmn flow] (cc) -- (as);
  \draw[bpmn flow] (as) |- (ch);
  \draw[bpmn flow] (as) |- (ci);
  \draw[bpmn flow] (ch) -| (aj);
  \draw[bpmn flow] (ci) -| (aj);
  \draw[bpmn flow] (aj) -- (aa);
  \draw[bpmn flow] (aa) -- (xs);
  \draw[bpmn flow] (xs) |- (mo) node[pos=0.85, above left] {granted (60\%)};
  \draw[bpmn flow] (xs) |- (nr) node[pos=0.85, below left] {denied (40\%)};
  \draw[bpmn flow] (mo) -| (xj);
  \draw[bpmn flow] (nr) -| (xj);
  \draw[bpmn flow] (xj) -- (end);
\end{tikzpicture}}
\end{center}
Decompose left to right and apply each rule: completeness check $1$\,d; AND-block $\max\{1,3\} = 3$\,d; assessment $3$\,d; XOR-block $0.6 \cdot 1 + 0.4 \cdot 2 = 1.4$\,d. The sequence rule sums the four fragments:
\[CT = 1 + \max\{1, 3\} + 3 + (0.6 \cdot 1 + 0.4 \cdot 2) = 8.4 \text{ days}.\]
\end{examplebox}

\paragraph{Rework loops.} Cyclic structures break the naïve sum: if $B\,(20)$ were always repeated, $CT = 10 + 20 + 20 + \cdots$ diverges. The missing ingredient is the \textbf{rework probability} $r$: the frequency with which the post-$B$ XOR-split sends the case back; with probability $1 - r$ the case exits. Two variants, ``$1$ or more'' (left) and ``$0$ or more'' (right):
\begin{center}
\begin{tikzpicture}[node distance=0.3cm and 0.4cm]
  \node[bpmn event] (start) {};
  \node[bpmn task, right=of start] (A) {$A\,(10)$};
  \node[bpmn gateway, right=of A] (xj) {$\times$};
  \node[bpmn task, right=of xj] (B) {$B\,(20)$};
  \node[bpmn gateway, right=of B] (xs) {$\times$};
  \node[bpmn event, draw=red, right=of xs] (end) {};
  \draw[bpmn flow] (start) -- (A);
  \draw[bpmn flow] (A) -- (xj);
  \draw[bpmn flow] (xj) -- (B);
  \draw[bpmn flow] (B) -- (xs);
  \draw[bpmn flow] (xs) -- (end) node[midway, above] {\scriptsize $1{-}r$};
  \draw[bpmn flow] (xs.south) -- ++(0,-0.45) node[pos=0.7, right] {\scriptsize $r$} -| (xj.south);
\end{tikzpicture}\hspace{0.8cm}
\begin{tikzpicture}[node distance=0.3cm and 0.4cm]
  \node[bpmn event] (start) {};
  \node[bpmn task, right=of start] (A) {$A\,(10)$};
  \node[bpmn gateway, right=of A] (xj) {$\times$};
  \node[bpmn gateway, right=1.9cm of xj] (xs) {$\times$};
  \node[bpmn task, below=0.45cm of $(xj)!0.5!(xs)$] (B) {$B\,(20)$};
  \node[bpmn event, draw=red, right=of xs] (end) {};
  \draw[bpmn flow] (start) -- (A);
  \draw[bpmn flow] (A) -- (xj);
  \draw[bpmn flow] (xj) -- (xs);
  \draw[bpmn flow] (xs) -- (end) node[midway, above] {\scriptsize $1{-}r$};
  \draw[bpmn flow] (xs.south) |- (B.east) node[pos=0.25, right] {\scriptsize $r$};
  \draw[bpmn flow] (B.west) -| (xj.south);
\end{tikzpicture}
\end{center}
In the left variant the body runs at least once; the case reworks $n - 1$ times and then exits with probability $r^{n-1}(1 - r)$, a geometric distribution over the iteration count. Weighting each count by its probability and writing $CT_P$ for the cycle time of the loop body:
\[CT = \sum_{i=1}^{\infty} i \cdot CT_P \cdot r^{i-1} (1 - r) = CT_P (1 - r) \sum_{i=1}^{\infty} i\, r^{i-1} = \frac{CT_P (1 - r)}{(1 - r)^2} = \frac{CT_P}{1 - r},\]
using the identity $\sum_{i \geq 1} i\,x^{i-1} = (1 - x)^{-2}$. The factor $1/(1 - r)$ inflates the body's time by the expected number of iterations.

In the right variant the case may skip $B$ entirely: exactly $n$ executions have probability $r^n (1 - r)$ for $n = 0, 1, 2, \ldots$, and the same computation gives
\[CT = \sum_{i=0}^{\infty} i \cdot CT_P \cdot r^{i} (1 - r) = CT_P (1 - r) \cdot \frac{r}{(1 - r)^2} = \frac{r \cdot CT_P}{1 - r}.\]
\begin{notebox}{Relating the two rules}
The $0$-or-more formula is the $1$-or-more one weighted by the probability $r$ that the body executes at all: a case skips $P$ with probability $1 - r$, contributing $0$, so $CT = \frac{CT_P}{1-r} \cdot r$.
\end{notebox}

\begin{examplebox}{Loan application with rework}
Refine the loan process: the completeness check ($2$\,h) is repeated until the application is complete, with rework probability $r = 0.2$ (the XOR-split after it exits with ``done'' $80\%$, loops back with ``repeat'' $20\%$); the parallel checks take $0.5$\,h (credit history) and $3$\,h (income sources); assessment $2$\,h; credit offer $2$\,h ($60\%$), rejection notice $0.5$\,h ($40\%$). The four patterns compose:
\[CT = \frac{2}{1 - 0.2} + \max\{0.5, 3\} + 2 + (0.6 \cdot 2 + 0.4 \cdot 0.5) = 2.5 + 3 + 2 + 1.4 = 8.9 \text{ hours}.\]
The rest composes as before.
\end{examplebox}

\begin{examplebox}{Inquiry handling: cycle time}
\begin{center}
\resizebox{0.95\linewidth}{!}{%
\begin{tikzpicture}[every node/.append style={font=\scriptsize\sffamily}]
  \node[bpmn event] (start) at (0,0) {};
  \node[bpmn task, align=center] (reg) at (1.6,0) {Register\\inquiry (2 d)};
  \node[bpmn task, align=center] (inv) at (4.1,0) {Investigate\\inquiry (8 d)};
  \node[bpmn gateway] (xj) at (5.9,0) {$\times$};
  \node[bpmn task, align=center] (pre) at (7.7,0) {Prepare\\response (4 d)};
  \node[bpmn task, align=center] (rev) at (10.3,0) {Review\\response (4 d)};
  \node[bpmn gateway] (xs) at (12.1,0) {$\times$};
  \node[bpmn event, draw=red] (end) at (13.3,0) {};
  \draw[bpmn flow] (start) -- (reg);
  \draw[bpmn flow] (reg) -- (inv);
  \draw[bpmn flow] (inv) -- (xj);
  \draw[bpmn flow] (xj) -- (pre);
  \draw[bpmn flow] (pre) -- (rev);
  \draw[bpmn flow] (rev) -- (xs);
  \draw[bpmn flow] (xs) -- (end) node[midway, above] {approved (80\%)};
  \draw[bpmn flow] (xs.south) -- ++(0,-0.7) node[pos=0.6, right] {rejected (20\%)} -| (xj.south);
\end{tikzpicture}}
\end{center}
The rework loop covers \emph{both} Prepare response and Review response: a rejected response is re-prepared and re-reviewed until approved, so $CT_P = 4 + 4 = 8$\,d with $r = 0.2$. The rest is a sequence:
\[CT = 2 + 8 + \frac{4 + 4}{1 - 0.2} = 10 + 10 = 20 \text{ days}.\]
The loop turns the $8$-day body into $10$ effective days: each case visits the prepare--review pair $1.25$ times on average.
\end{examplebox}

\subsection{Theoretical cycle time and efficiency}

In most real processes waiting time dominates the cycle time. Two causes recur: \emph{batched processing} (paperwork accumulated and handled in chunks) and \emph{actor unavailability} (approvals or prescriptions needing a specific person). Waiting time is therefore where redesign pays off most, and a pair of measures quantifies the opportunity.
\begin{definitionbox}{Theoretical cycle time, cycle time efficiency}
When the \emph{processing} time of every activity is known, the \textbf{theoretical cycle time} $TCT$ is computed by the same flow-analysis rules as $CT$ but with processing times in place of activity times: the cycle time the process would have if no case ever waited. The \textbf{cycle time efficiency} is
\[CTE = \frac{TCT}{CT} \in [0, 1].\]
$CTE \approx 1$ leaves little room to improve cycle time without radical changes; $CTE \approx 0$ signals a process dominated by idle time, improvable by cutting waits.
\end{definitionbox}
\begin{examplebox}{Inquiry handling: $TCT$ and $CTE$}
The inquiry process has $CT = 20$\,d $= 160$\,h (at $8$\,h per business day). Processing times: Register inquiry $30$\,min, Investigate inquiry $12$\,h, Prepare response $4$\,h, Review response $2$\,h. The same composition with processing times gives
\[TCT = 0.5 + 12 + \frac{4 + 2}{1 - 0.2} = 12.5 + 7.5 = 20 \text{ hours}, \qquad CTE = \frac{20}{160} = 12.5\%.\]
Seven eighths of the cycle time is waiting: a candidate for waiting-time reduction.
\end{examplebox}

\begin{notebox}{Limitations of flow analysis}
Three structural limits. The equations cover only \emph{block-structured} models (sequence, XOR-, AND-blocks, rework loops); goto-style jumps or unmatched gateways need ad-hoc reasoning. Activity times must be \emph{known}: estimating them, by interviews or logs, is its own measurement problem. And the formulas are \emph{load-insensitive}: they treat activity times as constants, while in reality waiting grows when arrivals increase against fixed resources.
\end{notebox}

\subsection{Little's law}

Two process-level measures accompany cycle time: the \textbf{arrival rate} $\lambda$, the average number of new cases created per unit of time, and the \textbf{Work-In-Process} $WIP$, the average number of cases active (started but not completed) at a given time. Both are averages over a stable window: $\lambda$ measures inflow, $WIP$ the active population.
\begin{theorem}[Little's law]
In a \emph{stable} system\footnote{Stable: the number of customers in the system does not grow unboundedly.} the long-term average number of customers equals the long-term average arrival rate times the average time a customer spends in the system:
\[WIP = \lambda \cdot CT.\]
\end{theorem}
John Little postulated it in 1954 and proved it in 1961; it holds with no assumption on arrival or service distributions and no knowledge of the queueing discipline.\footnote{J.~D.~C.~Little, \emph{A proof of the queuing formula $L = \lambda W$}, Operations Research 9 (1961), 383--387; the classical statement reads $L = \lambda \cdot W$.} Two readings matter. $WIP$ grows with either $CT$ or $\lambda$, so if $\lambda$ grows and $WIP$ must stay constant, $CT$ has to shrink. And since $WIP$ and $\lambda$ are cheap to sample, the law yields cycle time from two black-box observations, no flow analysis needed: $CT = WIP / \lambda$.
\begin{examplebox}{Cycle time by observation}
Over the past year ($250$ business days) $2500$ applications arrived: $\lambda = 2500/250 = 10$ per day. Weekly sampling found on average $200$ applications concurrently active: $WIP = 200$. Hence $CT = 200/10 = 20$ days: matching the flow-analysis result for the inquiry process.
\end{examplebox}
\begin{examplebox}{Restaurant}
A restaurant open $12$ hours serves $1200$ customers per day. During peak times ($6$ hours in total) it receives $900$ customers with on average $90$ present at any time; off-peak (the other $6$ hours) $300$ customers with $30$ present. Time spent by a customer in each regime:
\[CT_{\text{peak}} = \frac{90}{900/6} = 0.6 \text{ h}, \qquad CT_{\text{off}} = \frac{30}{300/6} = 0.6 \text{ h} = 36 \text{ minutes}.\]
The two regimes coincide: peak customers see a fuller room, but the inflow is proportionally higher. Now the manager expects peak demand to grow while the room is already at capacity. A higher $\lambda$ with $WIP$ capped by the building forces $CT$ down: customers must spend less time inside: shorten serving, redesign order-taking and payment, turn tables faster. Once physical capacity is fixed, process redesign is the only lever.
\end{examplebox}

\subsection{Cost analysis}

Flow analysis applies to cost with one change. Sequence, XOR-block and rework loops use the \emph{same} formulas as cycle time; the AND-block uses the \emph{sum} of the branch costs, not the maximum, because every branch incurs its cost regardless of which is slower: time runs in parallel, money does not. Activity cost itself has two components: \textbf{human-resource cost}, the hourly cost of the resource times the processing time, and \textbf{other cost}, fixed fees incurred per execution independently of the time spent.
\begin{examplebox}{Loan application: cost}
The refined loan process (rework on the completeness check, $r = 0.2$), with activities annotated by processing time. Clerk activities cost $25\,\euro$/h: Check completeness ($2$\,h), Check credit history ($0.5$\,h), Check income sources ($3$\,h). Credit-officer activities cost $50\,\euro$/h: Assess application ($2$\,h), Make credit offer ($2$\,h), Notify rejection ($0.5$\,h). Each credit-history check also costs a $1\,\euro$ fee.
\begin{center}
\begin{tabular}{llrrr}
\toprule
\textbf{Activity} & \textbf{Resource} & \textbf{Resource cost} & \textbf{Other} & \textbf{Total}\\
\midrule
Check completeness   & clerk   & $2 \cdot 25 = 50$     & $0$ & $50$\\
Check credit history & clerk   & $0.5 \cdot 25 = 12.5$ & $1$ & $13.5$\\
Check income sources & clerk   & $3 \cdot 25 = 75$     & $0$ & $75$\\
Assess application   & officer & $2 \cdot 50 = 100$    & $0$ & $100$\\
Make credit offer    & officer & $2 \cdot 50 = 100$    & $0$ & $100$\\
Notify rejection     & officer & $0.5 \cdot 50 = 25$   & $0$ & $25$\\
\bottomrule
\end{tabular}
\end{center}
Composing with the cost rules (rework: divide by $1 - r$; AND: sum; XOR: weighted average):
\[\text{Cost} = \frac{50}{1 - 0.2} + 13.5 + 75 + 100 + 0.6 \cdot 100 + 0.4 \cdot 25 = 62.5 + 188.5 + 70 = 321\,\euro.\]
The AND-block contributes $13.5 + 75 = 88.5\,\euro$ (the sum) where the cycle-time computation took only $\max\{0.5, 3\} = 3$ hours from the same block.
\end{examplebox}

\begin{examplebox}{Inquiry handling: cost}
The inquiry process again, with resources and rates: Register inquiry by a clerk ($25\,\euro$/h, $30$\,min), Investigate inquiry by an advisor ($50\,\euro$/h, $12$\,h), Prepare response by a senior advisor ($75\,\euro$/h, $4$\,h), Review response by a counselor ($100\,\euro$/h, $2$\,h). No AND-block; the rework loop covers the prepare--review pair with $r = 0.2$:
\[\text{Cost} = 25 \cdot 0.5 + 50 \cdot 12 + \frac{75 \cdot 4 + 100 \cdot 2}{1 - 0.2} = 12.5 + 600 + 625 = 1237.50\,\euro.\]
The loop sits on the costliest pair of activities and dominates the bill.
\end{examplebox}
